\thispagestyle{plain}
{\color{ink}\small\bfseries RESEARCH AT A GLANCE}\par
\vspace{2mm}
{\color{ink}\LARGE\bfseries Does more shopping detail produce better customer segments?}\par
\vspace{1mm}
{\color{ink}\rule{\textwidth}{0.8pt}}\par
\vspace{3mm}
\textbf{Finding.} Adding basket behavior and department preferences to RFM did not improve household segments. RFM K-means with two groups is the best-performing model: it is the most stable and separates later spending as well as any richer alternative.

\vfill
\begingroup
\setlength{\fboxsep}{10pt}
\colorbox{ink!7}{\parbox{\dimexpr\textwidth-20pt\relax}{
\centering\large\textbf{2.60 million transaction lines}\quad$\longrightarrow$\quad\textbf{2,498 households}\\[3mm]
\normalsize Days 1--500: build and freeze groups\quad$|$\quad Days 501--711: describe later behavior
}}
\endgroup

\vfill
{\large\bfseries\color{ink} Three feature spaces compared}\par
\vspace{2mm}
\begingroup
\setlength{\fboxsep}{0pt}
\newcommand{\spacecard}[4]{%
  \colorbox{ink!7}{\parbox[t][34mm][t]{\dimexpr(\textwidth-8mm)/3\relax}{%
    \vspace{3mm}\hspace*{3mm}\parbox[t][28mm][t]{\dimexpr(\textwidth-8mm)/3-6mm\relax}{%
      \raggedright{\color{ink}\large\bfseries #1\quad #2}\par\vspace{2mm}
      \small #3\par\vfill
      {\color{ink}\Large\bfseries #4}\ {\small features}}\vspace{3mm}}}}
\noindent\spacecard{01}{RFM}{Recency, frequency and monetary value}{3}\hfill
\spacecard{02}{Behavior}{RFM + basket size, variety and discount indicators}{8}\hfill
\spacecard{03}{Product}{Behavior + department spending shares}{17}
\endgroup

\vfill
{\large\bfseries\color{ink} Three results that change the recommendation}\par
\vspace{1mm}
{\renewcommand{\arraystretch}{1.3}
\begin{tabularx}{\textwidth}{@{}>{\raggedright\arraybackslash}p{36mm}>{\raggedright\arraybackslash}X>{\raggedright\arraybackslash}p{48mm}@{}}
\toprule
\textbf{Question} & \textbf{Evidence} & \textbf{Interpretation}\\\midrule
Will groups repeat? & Subsample ARI: \textbf{0.980 RFM} vs.\ 0.964 Product & RFM is the most stable.\\
Do groups separate later spend? & Explained variance: \textbf{20.4\% RFM} vs.\ 20.2\% Product & No gain from enrichment at $k=2$.\\
Does weighting matter? & At $k=3$, ARI to default falls to \textbf{0.433 / 0.460} at half / double category weight & Product segments depend on how similarity is defined.\\\bottomrule
\end{tabularx}}

\vfill
\begingroup
\setlength{\fboxsep}{8pt}
\fcolorbox{ink}{white}{\parbox{\dimexpr\textwidth-18pt\relax}{
\textbf{\color{ink}Decision.} Use RFM as the default for broad activity monitoring. Use department profiles as supplementary descriptions.
}}
\endgroup

\vspace{4mm}
\textbf{See the mechanism.} The household case studies show who changes groups; the membership-flow diagram makes weighting sensitivity visible. An exploratory later-category check finds a small Product advantage (0.386\% versus 0.299\% of share variation explained).

\vfill
{\small\color{ink}Note: the later-category extension was added after inspecting the original holdout results.}
\clearpage
